\documentclass[11pt]{article}
\usepackage[a4paper,margin=0.65in]{geometry}
\usepackage[T1]{fontenc}
\usepackage{amsmath,amssymb,amsfonts}
\usepackage{graphicx}
\usepackage{pdfpages}
\usepackage[english,shorthands=off]{babel}
\usepackage{fancyhdr}
\usepackage[bookmarks=false,colorlinks=false]{hyperref}
\usepackage[protrusion=true,expansion=false]{microtype}
\emergencystretch=3em
\tolerance=600
\providecommand{\modd}{\mathrm{mod}}
\providecommand{\Disc}{\mathrm{Disc}}
\providecommand{\canont}{}
\providecommand{\asterism}{*}
\providecommand{\Hessian}{H}
\providecommand{\tome}{}
\providecommand{\page}{p.}
\pagestyle{fancy}\fancyhf{}\fancyhead[L]{Pages 1--18}\fancyhead[R]{Cayley --- Collected Papers, Vol. IX}\renewcommand{\headrulewidth}{0.4pt}
\setlength{\headheight}{15pt}

\begin{document}

%% ========================================================================
%% Vol. IX front matter, PDF pages 1--18
%% Front cover, half-title, title page, advertisement, contents, blank.
%% No classification section (which begins on PDF p.~19).
%% No body papers (paper 556 begins on PDF p.~21).
%% ========================================================================

%% --- PDF pp. 1--4: outside cover, library stamps, blank flyleaves ---
\thispagestyle{empty}
\vspace*{\fill}
\begin{center}
\textit{[Front cover, library stamps (``FOR USE IN LIBRARY ONLY'', ``SEEN BY PRESERVATION SERVICES''), and blank flyleaves.]}
\end{center}
\vspace*{\fill}
\newpage

%% --- PDF p. 5: half-title ---
\thispagestyle{empty}
\vspace*{0.30\textheight}
\begin{center}
{\Large\scshape Mathematical Papers.}
\end{center}
\vspace*{\fill}
\newpage

%% --- PDF p. 6: publisher imprint ---
\thispagestyle{empty}
\vspace*{0.30\textheight}
\begin{center}
\textsl{London:} \textsc{C. J. Clay and Sons,}\\
\textsc{Cambridge University Press Warehouse,}\\
\textsc{Ave Maria Lane.}\\[2pt]
\textsl{Glasgow:} 263, \textsc{Argyle Street.}\\[10pt]
\textit{[Publisher's device.]}\\[10pt]
\textsl{Leipzig:} \textsc{F. A. Brockhaus.}\\
\textsl{New York:} \textsc{Macmillan and Co.}
\end{center}
\vspace*{\fill}
\newpage

%% --- PDF p. 7: title page ---
\thispagestyle{empty}
\begin{center}
\vspace*{0.06\textheight}
{\LARGE\scshape The Collected}\\[1.2em]
{\Huge\scshape Mathematical Papers}\\[1.0em]
{\large\scshape of}\\[1.0em]
{\Large\scshape Arthur Cayley, Sc.D., F.R.S.,}\\[0.4em]
{\footnotesize\scshape late Sadlerian Professor of Pure Mathematics in the University of Cambridge.}\\[3.2em]
{\Large\scshape Vol. IX.}\\[3.2em]
{\Large\scshape Cambridge:}\\
{\large\scshape At the University Press.}\\[0.3em]
{\large 1896}\\[1.4em]
{\footnotesize \textit{[All Rights reserved.]}}
\end{center}
\vspace*{\fill}
\newpage

%% --- PDF p. 8: printer's imprint (verso of title) ---
\thispagestyle{empty}
\vspace*{0.45\textheight}
\begin{center}
\textsc{Cambridge:}\\
\textsc{Printed by J. and C. F. Clay,}\\
\textsc{At the University Press.}
\end{center}
\vspace*{\fill}
\newpage

%% --- PDF p. 9: Advertisement (Roman v) ---
\thispagestyle{empty}
\begin{flushright}\small v\end{flushright}
\vspace*{1.5em}
\begin{center}
{\Large\scshape Advertisement.}
\end{center}
\vspace*{1em}

\noindent T\textsc{he} present volume contains 74 papers, numbered 556 to 629, published for the most part in the years 1874 to 1877.

\medskip
The Table for the nine volumes is
\begin{center}
\begin{tabular}{rlrrl}
Vol. & I.    & Numbers \quad 1   & to & 100,\\
     & II.   & ,,\quad\phantom{0}101 & ,, & 158,\\
     & III.  & ,,\quad\phantom{0}159 & ,, & 222,\\
     & IV.   & ,,\quad\phantom{0}223 & ,, & 299,\\
     & V.    & ,,\quad\phantom{0}300 & ,, & 383,\\
     & VI.   & ,,\quad\phantom{0}384 & ,, & 416,\\
     & VII.  & ,,\quad\phantom{0}417 & ,, & 485,\\
     & VIII. & ,,\quad\phantom{0}486 & ,, & 555,\\
     & IX.   & ,,\quad\phantom{0}556 & ,, & 629.
\end{tabular}
\end{center}

\vspace{2em}
\hfill A.~R.~F\textsc{orsyth}.

\medskip
\noindent\textit{17 December 1895.}
\vspace*{\fill}
\newpage

%% --- PDF p. 10: Internet Archive digitization notice ---
\thispagestyle{empty}
\vspace*{\fill}
\begin{center}
\textit{[Digitized by the Internet Archive in 2007 with funding from Microsoft Corporation. URL: \texttt{http://www.archive.org/details/collectedmathema09cayluoft}]}
\end{center}
\vspace*{\fill}
\newpage

%% --- PDF p. 11: Contents (vii) ---
\thispagestyle{empty}
\begin{flushright}\small vii\end{flushright}
\vspace*{1em}
\begin{center}
{\Large\scshape Contents.}
\end{center}
\vspace*{1em}

\noindent\hfill\textsc{page}

\medskip

\begin{description}\itemsep=2pt
\item[556.] \textit{On Steiner's surface}\dotfill 1\\
\hspace*{1em}\small Proc. Lond. Math. Society, t.~v.\ (1873--1874), pp.~14--25.

\item[557.] \textit{On certain constructions for bicircular quartics}\dotfill 13\\
\hspace*{1em}\small Proc. Lond. Math. Society, t.~v.\ (1873--1874), pp.~29--31.

\item[558.] \textit{A geometrical interpretation of the equations obtained by equating to zero the resultant and the discriminants of two binary quantics}\dotfill 16\\
\hspace*{1em}\small Proc. Lond. Math. Society, t.~v.\ (1873--1874), pp.~31--33.

\item[559.] \textit{[Note on inversion]}\dotfill 18\\
\hspace*{1em}\small Proc. Lond. Math. Society, t.~v.\ (1873--1874), p.~112.

\item[560.] \textit{[Addition to Lord Rayleigh's paper ``On the numerical calculation of the roots of fluctuating functions'']}\dotfill 19\\
\hspace*{1em}\small Proc. Lond. Math. Society, t.~v.\ (1873--1874), pp.~123, 124.

\item[561.] \textit{On the geometrical representation of Cauchy's theorems of root-limitation}\dotfill 21\\
\hspace*{1em}\small Camb.\ Phil.\ Trans., t.~xii.\ Part~ii.\ (1877), pp.~395--413.

\item[562.] \textit{On a theorem in maxima and minima: addition [to Mr Walton's paper] by Professor Cayley}\dotfill 40\\
\hspace*{1em}\small Quart. Math. Jour., t.~x.\ (1870), pp.~262, 263.

\item[563.] \textit{Note on the transformation of two simultaneous equations}\dotfill 42\\
\hspace*{1em}\small Quart. Math. Jour., t.~xi.\ (1871), pp.~266, 267.

\item[564.] \textit{On a theorem in elimination}\dotfill 43\\
\hspace*{1em}\small Quart. Math. Jour., t.~xii.\ (1873), pp.~5, 6.
\end{description}
\newpage

%% --- PDF p. 12: Contents (viii) ---
\thispagestyle{empty}
\noindent\small viii\hfill\textsc{contents.}
\vspace*{1em}

\noindent\hfill\textsc{page}

\medskip

\begin{description}\itemsep=2pt
\item[565.] \textit{Note on the Cartesian}\dotfill 45\\
\hspace*{1em}\small Quart. Math. Jour., t.~xii.\ (1873), pp.~16--19.

\item[566.] \textit{On the transformation of the equation of a surface to a set of chief axes}\dotfill 48\\
\hspace*{1em}\small Quart. Math. Jour., t.~xii.\ (1873), pp.~34--38.

\item[567.] \textit{On an identical equation connected with the theory of invariants}\dotfill 52\\
\hspace*{1em}\small Quart. Math. Jour., t.~xii.\ (1873), pp.~115--118.

\item[568.] \textit{Note on the integrals $\int_{0}^{\infty}\cos x^{2}\,dx$ and $\int_{0}^{\infty}\sin x^{2}\,dx$}\dotfill 56\\
\hspace*{1em}\small Quart. Math. Jour., t.~xii.\ (1873), pp.~118--126.

\item[569.] \textit{On the cyclide}\dotfill 64\\
\hspace*{1em}\small Quart. Math. Jour., t.~xii.\ (1873), pp.~148--165.

\item[570.] \textit{On the superlines of a quadric surface in five-dimensional space}\dotfill 79\\
\hspace*{1em}\small Quart. Math. Jour., t.~xii.\ (1873), pp.~176--180.

\item[571.] \textit{A demonstration of Dupin's theorem}\dotfill 84\\
\hspace*{1em}\small Quart. Math. Jour., t.~xii.\ (1873), pp.~185--191.

\item[572.] \textit{Theorem in regard to the Hessian of a quaternary function}\dotfill 90\\
\hspace*{1em}\small Quart. Math. Jour., t.~xii.\ (1873), pp.~193--197.

\item[573.] \textit{Note on the $(2,2)$ correspondence of two variables}\dotfill 94\\
\hspace*{1em}\small Quart. Math. Jour., t.~xii.\ (1873), pp.~197, 198.

\item[574.] \textit{On Wronski's theorem}\dotfill 96\\
\hspace*{1em}\small Quart. Math. Jour., t.~xii.\ (1873), pp.~221--228.

\item[575.] \textit{On a special quartic transformation of an elliptic function}\dotfill 103\\
\hspace*{1em}\small Quart. Math. Jour., t.~xii.\ (1873), pp.~266--269.

\item[576.] \textit{Addition to Mr Walton's paper ``On the ray-planes in biaxal crystals''}\dotfill 107\\
\hspace*{1em}\small Quart. Math. Jour., t.~xii.\ (1873), pp.~273--275.

\item[577.] \textit{Note in illustration of certain general theorems obtained by Dr Lipschitz}\dotfill 110\\
\hspace*{1em}\small Quart. Math. Jour., t.~xii.\ (1873), pp.~346--349.
\end{description}
\newpage

%% --- PDF p. 13: Contents (ix) ---
\thispagestyle{empty}
\noindent\small\textsc{contents.}\hfill ix
\vspace*{1em}

\noindent\hfill\textsc{page}

\medskip

\begin{description}\itemsep=2pt
\item[578.] \textit{A memoir on the transformation of elliptic functions}\dotfill 113\\
\hspace*{1em}\small Phil. Trans., t.~clxiv.\ (for 1874), pp.~397--456.

\item[579.] \textit{Address delivered by the President, Professor Cayley, on presenting the Gold Medal of the [Royal Astronomical] Society to Professor Simon Newcomb}\dotfill 176\\
\hspace*{1em}\small Monthly Notices R.~Ast.\ Society, t.~xxxiv.\ (1873--1874), pp.~224--233.

\item[580.] \textit{On the number of distinct terms in a symmetrical or partially symmetrical determinant; with an addition}\dotfill 185\\
\hspace*{1em}\small Monthly Notices R.~Ast.\ Society, t.~xxxiv.\ (1873--1874), pp.~303--307; p.~335.

\item[581.] \textit{On a theorem in elliptic motion}\dotfill 191\\
\hspace*{1em}\small Monthly Notices R.~Ast.\ Society, t.~xxxv.\ (1874--1875), pp.~337--339.

\item[582.] \textit{Note on the Theory of Precession and Nutation}\dotfill 194\\
\hspace*{1em}\small Monthly Notices R.~Ast.\ Society, t.~xxxv.\ (1874--1875), pp.~340--343.

\item[583.] \textit{On spheroidal trigonometry}\dotfill 197\\
\hspace*{1em}\small Monthly Notices R.~Ast.\ Society, t.~xxxvii.\ (1876--1877), p.~92.

\item[584.] \textit{Addition to Prof.\ R.~S.\ Ball's paper ``Note on a transformation of Lagrange's equations of motion in generalised coordinates, which is convenient in Physical Astronomy''}\dotfill 198\\
\hspace*{1em}\small Monthly Notices R.~Ast.\ Society, t.~xxxvii.\ (1876--1877), pp.~269--271.

\item[585.] \textit{A new theorem on the equilibrium of four forces acting on a solid body}\dotfill 201\\
\hspace*{1em}\small Phil. Mag., t.~xxxi.\ (1866), pp.~78, 79; Camb.\ Phil.\ Soc.\ Proc., t.~i.\ (1866), p.~235.

\item[586.] \textit{On the mathematical theory of isomers}\dotfill 202\\
\hspace*{1em}\small Phil. Mag., t.~xlvii.\ (1874), pp.~444--467.

\item[587.] \textit{A Smith's Prize dissertation [1873]}\dotfill 205\\
\hspace*{1em}\small Messenger of Mathematics, t.~iii.\ (1874), pp.~1--4.

\item[588.] \textit{Problem [on tetrahedra]}\dotfill 209\\
\hspace*{1em}\small Messenger of Mathematics, t.~iii.\ (1874), pp.~50--52.
\end{description}

\vspace{1em}
\noindent\small C.~IX.\hfill b
\newpage

%% --- PDF p. 14: Contents (x) ---
\thispagestyle{empty}
\noindent\small x\hfill\textsc{contents.}
\vspace*{1em}

\noindent\hfill\textsc{page}

\medskip

\begin{description}\itemsep=2pt
\item[589.] \textit{On residuation in regard to a cubic curve}\dotfill 211\\
\hspace*{1em}\small Messenger of Mathematics, t.~iii.\ (1874), pp.~62--65.

\item[590.] \textit{Addition to Prof.\ Hall's paper ``On the motion of a particle toward an attracting centre at which the force is infinite.''}\dotfill 215\\
\hspace*{1em}\small Messenger of Mathematics, t.~iii.\ (1874), pp.~149--152.

\item[591.] \textit{A Smith's Prize paper and dissertation [1874]; solutions and remarks}\dotfill 218\\
\hspace*{1em}\small Messenger of Mathematics, t.~iii.\ (1874), pp.~165--183; t.~iv.\ (1875), pp.~6--8.

\item[592.] \textit{On the Mercator's projection of a skew hyperboloid of revolution}\dotfill 237\\
\hspace*{1em}\small Messenger of Mathematics, t.~iv.\ (1875), pp.~17--20.

\item[593.] \textit{A Sheepshanks' problem (1866)}\dotfill 241\\
\hspace*{1em}\small Messenger of Mathematics, t.~iv.\ (1875), pp.~34--36.

\item[594.] \textit{On a differential equation in the theory of elliptic functions}\dotfill 244\\
\hspace*{1em}\small Messenger of Mathematics, t.~iv.\ (1875), pp.~69, 70.

\item[595.] \textit{On a Senate-House problem}\dotfill 246\\
\hspace*{1em}\small Messenger of Mathematics, t.~iv.\ (1875), pp.~75--78.

\item[596.] \textit{Note on a theorem of Jacobi's for the transformation of a double integral}\dotfill 250\\
\hspace*{1em}\small Messenger of Mathematics, t.~iv.\ (1875), pp.~92--94.

\item[597.] \textit{On a differential equation in the theory of elliptic functions}\dotfill 253\\
\hspace*{1em}\small Messenger of Mathematics, t.~iv.\ (1875), pp.~110--113.

\item[598.] \textit{Note on a process of integration}\dotfill 257\\
\hspace*{1em}\small Messenger of Mathematics, t.~iv.\ (1875), pp.~149, 150.

\item[599.] \textit{A Smith's Prize dissertation}\dotfill 259\\
\hspace*{1em}\small Messenger of Mathematics, t.~iv.\ (1875), pp.~157--160.

\item[600.] \textit{Theorem on the $n$-th Roots of Unity}\dotfill 263\\
\hspace*{1em}\small Messenger of Mathematics, t.~iv.\ (1875), p.~171.

\item[601.] \textit{Note on the Cassinian}\dotfill 264\\
\hspace*{1em}\small Messenger of Mathematics, t.~iv.\ (1875), pp.~187, 188.
\end{description}
\newpage

%% --- PDF p. 15: Contents (xi) ---
\thispagestyle{empty}
\noindent\small\textsc{contents.}\hfill xi
\vspace*{1em}

\noindent\hfill\textsc{page}

\medskip

\begin{description}\itemsep=2pt
\item[602.] \textit{On the potentials of polygons and polyhedra}\dotfill 266\\
\hspace*{1em}\small Proc. Lond. Math. Society, t.~vi.\ (1874--1875), pp.~20--34.

\item[603.] \textit{On the potential of the ellipse and the circle}\dotfill 281\\
\hspace*{1em}\small Proc. Lond. Math. Society, t.~vi.\ (1874--1875), pp.~38--58.

\item[604.] \textit{Determination of the attraction of an ellipsoidal shell on an exterior point}\dotfill 302\\
\hspace*{1em}\small Proc. Lond. Math. Society, t.~vi.\ (1874--1875), pp.~58--67.

\item[605.] \textit{Note on a point in the theory of attraction}\dotfill 312\\
\hspace*{1em}\small Proc. Lond. Math. Society, t.~vi.\ (1874--1875), pp.~79--81.

\item[606.] \textit{On the expression of the coordinates of a point of a quartic curve as functions of a parameter}\dotfill 315\\
\hspace*{1em}\small Proc. Lond. Math. Society, t.~vi.\ (1874--1875), pp.~81--83.

\item[607.] \textit{A memoir on prepotentials}\dotfill 318\\
\hspace*{1em}\small Phil. Trans., t.~clxv.\ (for 1875), pp.~675--774.

\item[608.] \textit{[Extract from a] Report on Mathematical Tables}\dotfill 424\\
\hspace*{1em}\small Brit. Assoc.\ Report, 1873, pp.~3, 4.

\item[609.] \textit{On the analytical forms called factions}\dotfill 426\\
\hspace*{1em}\small Brit. Assoc.\ Report, 1875, Notices of Communications to the Sections, p.~10.

\item[610.] \textit{On the analytical forms called Trees, with application to the theory of chemical combinations}\dotfill 427\\
\hspace*{1em}\small Brit. Assoc.\ Report, 1875, pp.~257--305.

\item[611.] \textit{Report on mathematical tables}\dotfill 461\\
\hspace*{1em}\small Brit. Assoc.\ Report, 1875, pp.~305--336.

\item[612.] \textit{Note sur une formule d'int\'egration ind\'efinie}\dotfill 500\\
\hspace*{1em}\small Comptes Rendus, t.~lxxviii.\ (1874), pp.~1624--1629.

\item[613.] \textit{On the group of points $G_4^1$ on a sextic curve with five double points}\dotfill 504\\
\hspace*{1em}\small Math. Ann., t.~viii.\ (1875), pp.~359--362.

\item[614.] \textit{On a problem of projection}\dotfill 508\\
\hspace*{1em}\small Quart. Math. Jour., t.~xiii.\ (1875), pp.~19--29.
\end{description}

\vspace{1em}
\hfill\small b 2
\newpage

%% --- PDF p. 16: Contents (xii) ---
\thispagestyle{empty}
\noindent\small xii\hfill\textsc{contents.}
\vspace*{1em}

\noindent\hfill\textsc{page}

\medskip

\begin{description}\itemsep=2pt
\item[615.] \textit{On the conic torus}\dotfill 519\\
\hspace*{1em}\small Quart. Math. Jour., t.~xiii.\ (1875), pp.~127--129.

\item[616.] \textit{A geometrical illustration of the cubic transformation in elliptic functions}\dotfill 522\\
\hspace*{1em}\small Quart. Math. Jour., t.~xiii.\ (1875), pp.~211--216.

\item[617.] \textit{On the scalene transformation of a plane curve}\dotfill 527\\
\hspace*{1em}\small Quart. Math. Jour., t.~xiii.\ (1875), pp.~321--328.

\item[618.] \textit{On the mechanical description of a Cartesian}\dotfill 535\\
\hspace*{1em}\small Quart. Math. Jour., t.~xiii.\ (1875), pp.~328--330.

\item[619.] \textit{On an algebraical operation}\dotfill 537\\
\hspace*{1em}\small Quart. Math. Jour., t.~xiii.\ (1875), pp.~369--375.

\item[620.] \textit{Correction of two numerical errors in Sohnke's paper respecting modular equations}\dotfill 543\\
\hspace*{1em}\small Crelle, t.~lxxxi.\ (1876), p.~229.

\item[621.] \textit{On the number of the univalent radicals $\mathrm{C}_{n}\mathrm{H}_{2n+1}$}\dotfill 544\\
\hspace*{1em}\small Phil. Mag., Ser.~5, t.~iii.\ (1877), pp.~34, 35.

\item[622.] \textit{On a system of equations connected with Malfatti's problem}\dotfill 546\\
\hspace*{1em}\small Proc. Lond. Math. Society, t.~vii.\ (1876), pp.~38--42.

\item[623.] \textit{On three-bar motion}\dotfill 551\\
\hspace*{1em}\small Proc. Lond. Math. Society, t.~vii.\ (1876), pp.~136--166.

\item[624.] \textit{On the bicursal sextic}\dotfill 581\\
\hspace*{1em}\small Proc. Lond. Math. Society, t.~vii.\ (1876), pp.~166--172.

\item[625.] \textit{On the condition for the existence of a surface cutting at right angles a given set of lines}\dotfill 587\\
\hspace*{1em}\small Proc. Lond. Math. Society, t.~viii.\ (1877), pp.~53--57.

\item[626.] \textit{On the general differential equation $\dfrac{dx}{\sqrt{X}}+\dfrac{dy}{\sqrt{Y}}=0$, where $X$, $Y$ are the same quartic functions of $x$, $y$ respectively}\dotfill 592\\
\hspace*{1em}\small Proc. Lond. Math. Society, t.~viii.\ (1877), pp.~184--199.
\end{description}
\newpage

%% --- PDF p. 17: Contents (xiii), end of contents ---
\thispagestyle{empty}
\noindent\small\textsc{contents.}\hfill xiii
\vspace*{1em}

\noindent\hfill\textsc{page}

\medskip

\begin{description}\itemsep=2pt
\item[627.] \textit{Geometrical illustration of a theorem relating to an irrational function of an imaginary variable}\dotfill 609\\
\hspace*{1em}\small Proc. Lond. Math. Society, t.~viii.\ (1877), pp.~212--214.

\item[628.] \textit{On the circular relation of M\"obius}\dotfill 612\\
\hspace*{1em}\small Proc. Lond. Math. Society, t.~viii.\ (1877), pp.~220--225.

\item[629.] \textit{On the linear transformation of the integral $\displaystyle\int\frac{du}{\sqrt{U}}$}\dotfill 618\\
\hspace*{1em}\small Proc. Lond. Math. Society, t.~viii.\ (1877), pp.~226--229.
\end{description}

\vspace{2em}
\noindent Plate \dotfill \textit{to face p.\ 460}
\newpage

%% --- PDF p. 18: blank ---
\thispagestyle{empty}
\vspace*{\fill}
\begin{center}
\textit{[Blank page.]}
\end{center}
\vspace*{\fill}

\end{document}
